% Generated by paper/tables.py -- do not edit by hand.
% Table 15: anchor transition temperatures, surrogate vs CALPHAD vs DTA
\footnotesize
\setlength{\tabcolsep}{4pt}
\resizebox{\textwidth}{!}{
\begin{tabular}{llrrrrrr}
\toprule
& & \multicolumn{3}{c}{transition temperature ($^{\circ}$C)} & \multicolumn{3}{c}{$\Delta T$ (K)} \\
\cmidrule(lr){3-5}\cmidrule(lr){6-8}
Alloy & Composition (wt\%) & Published & CALPHAD & Surrogate & sur$-$cal & cal$-$exp & sur$-$exp \\
\midrule
Drozdov\'a A & C 0.308, Cr 1.058 & $T_{\mathrm{L}}$ 1498 & 1511 & 1525 & +14 & +13 & +27 \\
Drozdov\'a B & C 0.32, Cr 1.54 & $T_{\mathrm{P}}$ 1471 & 1469 & 1450 & -19 & -2 & -21 \\
Drozdov\'a C & C 0.38, Cr 4.99 & $T_{\mathrm{S}}$ 1397 & 1456 & 1435 & -21 & +59 & +38 \\
Yamada 21Cr--16Ni & Cr 21.0, Ni 16.6 & $T_{\mathrm{L}}$ 1441 & 1438 & 1455 & +17 & -3 & +14 \\
\midrule
\multicolumn{8}{l}{\scriptsize $T_L\leftrightarrow$ last solid, $T_P$/$T_S\leftrightarrow$ first liquid; detected at liquid fraction $10^{-4}$ on the sustained branch.} \\
\multicolumn{8}{l}{\scriptsize Composition in wt\% (Fe balance); surrogate = 3-seed ensemble, primary head (Fe--Cr--C: sigmoid/$\Sigma$; Fe--Cr--Ni: renorm).} \\
\bottomrule
\end{tabular}}
